\section{Setup}
\label{sec:setup}

\subsection{Model and data}

All experiments use \texttt{Qwen/Qwen3-4B-Instruct-2507} \citep{qwen3report2025}, a 36-layer decoder-only transformer with a 2560-dimensional hidden state. The dataset consists of 444 LiveCodeBench \citep{jain2024livecodebench} tasks spanning the easy, medium, and hard difficulty tiers. Every experiment runs on a single NVIDIA DGX system.

\subsection{Hidden state capture}

Across all task attempts, the hidden state at the final prompt token is extracted from a prompt-only forward pass, before the autoregressive continuation begins. This yields a single 2560-dimensional vector per task per attempt at each captured layer. Hidden states are captured for the upper transformer blocks (layers 29--36), motivated by the general finding that linear separability of high-level properties increases with depth in deep networks \citep{alain2016probes}. Layer selection is handled by nested cross-validation (CV) in Section~\ref{sec:probe_results}.

\subsection{Repair-B policy}

For tasks where the model's first attempt fails the unit tests, an iterative repair procedure (``repair-B'') is invoked, with a maximum of 5 total attempts per task. Each attempt is a fresh forward pass on the full prompt, with no state carried over from previous attempts, and the repair prompt contains the original problem, the model's previous code, the error output captured by the verifier, and an instruction to fix the failure. Generation is capped at 1024 new tokens per attempt. The pre-generation probe in Section~\ref{sec:probe_results} uses all 444 tasks at attempt 0, and the outcome of the repair procedure itself is summarized in Section~\ref{sec:limitations}.
